\documentclass[uplatex,dvipdfmx]{jsarticle}

\usepackage{amsmath,amssymb,amsthm}
\usepackage{mathtools}

\newtheorem{theorem}{定理}
\newtheorem{proposition}{命題}
\newtheorem{lemma}{補題}
\theoremstyle{definition}
\newtheorem{definition}{定義}
\theoremstyle{remark}
\newtheorem{remark}{注意}

\begin{document}

\section*{正則値の逆像の接空間}

\begin{proposition}
$M$ を $n$ 次元滑らかな多様体とし，
$P$ を $m$ 次元滑らかな多様体とする。
$F:M\to P$ を滑らかな写像とする。
$c\in P$ が $F$ の正則値であるとし，
\[
  N=F^{-1}(c)
\]
とおく。
ただし $N\ne\varnothing$ と仮定する。

このとき $N$ は $M$ の余次元 $m$ の埋め込み部分多様体であり，
包含写像を
\[
  i:N\hookrightarrow M
\]
と書くと，任意の $p\in N$ に対して
\[
  i_{*,p}(T_pN)=\ker F_{*,p}\subset T_pM
\]
が成り立つ。

したがって，$i_{*,p}$ によって $T_pN$ を $T_pM$ の部分空間と同一視すれば，
\[
  T_pN=\ker F_{*,p}
\]
と書ける。
\end{proposition}

\begin{proof}
まず，$c$ が $F$ の正則値であるから，任意の $p\in N=F^{-1}(c)$ に対して
\[
  F_{*,p}:T_pM\to T_cP
\]
は全射である。

特に，正則値定理により，$N=F^{-1}(c)$ は $M$ の余次元 $m$ の埋め込み部分多様体である。
以下，任意に $p\in N$ を固定する。

まず
\[
  i_{*,p}(T_pN)\subset \ker F_{*,p}
\]
を示す。

$N=F^{-1}(c)$ なので，$F$ は $N$ 上で定数写像 $c$ に等しい。
したがって合成写像
\[
  F\circ i:N\to P
\]
は定数写像である。ゆえに任意の $\xi\in T_pN$ に対して
\[
  (F\circ i)_{*,p}(\xi)=0
\]
である。

一方，接写像の連鎖律より
\[
  (F\circ i)_{*,p}=F_{*,p}\circ i_{*,p}
\]
であるから，
\[
  F_{*,p}\bigl(i_{*,p}(\xi)\bigr)=0
\]
となる。したがって
\[
  i_{*,p}(\xi)\in \ker F_{*,p}
\]
である。よって
\[
  i_{*,p}(T_pN)\subset \ker F_{*,p}
\]
が従う。

あとは次元を比較する。
正則値定理により $N$ は $M$ の余次元 $m$ の埋め込み部分多様体であるから，
\[
  \dim T_pN=n-m
\]
である。
また，包含写像 $i:N\hookrightarrow M$ は埋め込みであるから，
$i_{*,p}:T_pN\to T_pM$ は単射である。
したがって
\[
  \dim i_{*,p}(T_pN)=n-m
\]
である。

一方，$F_{*,p}:T_pM\to T_cP$ は全射であり，
$\dim T_pM=n$，$\dim T_cP=m$ である。
よって階数・退化次数の定理より
\[
  \dim \ker F_{*,p}=n-m
\]
である。

すでに
\[
  i_{*,p}(T_pN)\subset \ker F_{*,p}
\]
を示しており，両辺はともに $T_pM$ の $n-m$ 次元部分空間である。
したがって
\[
  i_{*,p}(T_pN)=\ker F_{*,p}
\]
である。
\end{proof}

\begin{remark}
慣習的には，包含写像
\[
  i:N\hookrightarrow M
\]
の接写像
\[
  i_{*,p}:T_pN\to T_pM
\]
によって $T_pN$ を $T_pM$ の部分空間とみなす。
この同一視のもとで，上の結論は簡潔に
\[
  T_pN=\ker F_{*,p}
\]
と書かれる。
\end{remark}

\end{document}
